\section{Complete reference C from the first homogeneous rank}
Let $H=(1-z-z^2)^{-1}$ and obtain $F_V$ by repeatedly applying
$F_{Vt}=1+S_tF_V$, where $S_t=z+\cdots+z^t$ and $t\ge2$.
Remove all leading twos, which fix $H$. For the remaining nonempty
word $V$ put $W=\sum t$, $r=|V|$, $B=(1-z-z^2)F_V$,
$d=(1-z)F_V$ and $G=F_V/(1-z)$. The word need not be sorted.
The existing numerator argument proves $B\ge0$, $B_0=1$,
$\deg B=W$, and $d_n>0$ for $n\ge2$.

\begin{theorem}[All shifts at every homogeneous rank]
For EVERY finite word as above, every $N\ge W+1$ and every $R\ge3$,
\[
                    G_{N-R}d_{N-1}\ge G_{N-R-1}d_N.       \tag{1}
\]
Negative indices are zero. Empty effective words have the already
proved complete binary C. This improves the previous universal
start $W+6$ to $W+1$. It does not assert C at every earlier rank.
\end{theorem}

\subsection*{A lower ratio envelope using the binary cumulative sequence}
Write $f_n=\operatorname{Fib}(n)$ and $g_k=f_{k+3}-1$ for $k\ge0$,
with $g_k=0$ for $k<0$. This is the cumulative sequence of $H$.
For $k\ge1$, its LC minor is exactly
\[
                   g_k^2-g_{k-1}g_{k+1}=f_k+(-1)^k\ge0. \tag{2}
\]
Indeed the Fibonacci Cassini term is $(-1)^k$, and the remaining
linear term is $f_{k+2}+f_{k+4}-2f_{k+3}=f_k$.
Thus for $1\le k\le n$ and every $i\ge0$,
\[
                 g_{k-i}g_{n-1}\ge g_{k-i-1}g_n.
\]
For negative lower indices this follows by zero extension; otherwise
it is an ordinary ratio comparison in the positive LC sequence $g$.
Multiplication by $B_i$ and summation give the all-read envelope
\[
                  G_k g_{n-1}\ge G_{k-1}g_n
                            \qquad(1\le k\le n).        \tag{3}
\]
No LC of the arbitrary-reference cumulative sequence $G$ is assumed.
Consequently it suffices for all the comparisons at rank $N\ge4$
to prove the SINGLE comparison
\[
           (f_N-1)d_{N-1}\ge(f_{N-1}-1)d_N.              \tag{4}
\]
The reads $k\le0$ are automatic.

\subsection*{A relative spectral bound for the entire word}
Let $p=\varphi$, $q=p^{-1}=p-1$ and $\psi=-q$.
Evaluation at either zero of $1-z-z^2$ gives
\[
 A:=B(q)=\prod_{t\in V}S_t(q)>0,\qquad
 C:=B(-p)=\prod_{t\in V}S_t(-p).
\]
The exact all-capacity identities are
\[
 pS_t(-p)=(-p)^t-1,\qquad
 p^t S_t(q)=p(p^t-1).                                  \tag{5}
\]
For EVERY $t\ge2$ they imply
\[
                    p|S_t(-p)|\le p^t S_t(q).           \tag{6}
\]
For $t=2$ both windows are one. For $t\ge3$, the triangle
inequality and $p^t\ge p^3$ give
\[
 p|S_t(-p)|\le p^t+1\le p(p^t-1).
\]
The last inequality is $(p-1)p^t\ge p+1=p^2$ and is an equality
at $t=3$. Multiplying (6) over the WHOLE word gives
\[
                         p^r|C|\le p^W A.              \tag{7}
\]
This relative bound retains the positive-root amplitude $A$;
the older absolute bound $|C|\le p^W$ loses precisely this resource.
Also $\operatorname{sgn}(C)=(-1)^W$, since the first identity in (5)
shows $\operatorname{sgn}(S_t(-p))=(-1)^t$.

\subsection*{The exact comparison at all homogeneous ranks}
The binary signed slope at rank $j\ge0$ is $f_{j-1}$, using
$f_{-1}=1$. The exceptional atom at zero is therefore INCLUDED.
Binet's formula consequently gives, already for $n\ge W$,
\[
               \sqrt5\,d_n=A p^{n-1}-C\psi^{n-1}.       \tag{8}
\]
Thus at $N\ge W+1$ the conjugate error is exactly
\[
                     d_N-pd_{N-1}=C\psi^{N-2}.          \tag{9}
\]
If $N-W$ is odd, (9) is negative. Since $g_k-pg_{k-1}>0$
for every nonnegative $k$, (1) follows directly. In particular
the FIRST homogeneous rank $N=W+1$ is always paid, with no norm bound.

Suppose now $N\ge W+2$ and $r\ge2$. By (7),
\[
                              |C|\le A p^{N-4}.         \tag{10}
\]
Put $\sigma=(-1)^{N-2}$. Using (8), the exact determinant in (4) is
\[
\begin{split}
 \sqrt5\{(f_N-1)d_{N-1}-(f_{N-1}-1)d_N\}
  =A(p^{N-3}-\sigma p^{-1})
                  -\sigma C p(1-p^{-N+2}).             \tag{11}
\end{split}
If $\sigma C\le0$, both terms have nonnegative total sign directly.
Otherwise (10) bounds (11) below by
\[
 A\{p^{N-3}-\sigma p^{-1}-p^{N-3}(1-p^{-N+2})\}
                         =A p^{-1}(1-\sigma)\ge0.
\]
This proves (4). Applying (3) with $n=N-3$ proves EVERY $R\ge3$.
For $r=1$, write $V=(t)$. At $N=W+2=t+2$ the explicit one-window
slopes give $d_{N-1}=f_{N-1}-1$ and $d_N=f_N-1$, so (4) is
an equality. For $N\ge W+3$, (7) again implies (10), since
$W-r=W-1\le N-4$, and (11) proves (4). The first rank $W+1$
was already paid by its negative conjugate error. Thus the proof
does not require an imported early one-jump C theorem.
This completes the proof of (1).

\paragraph{An exact equality checking the delicate boundary.}
For $V=(3,3)$, $W=6$, the reference slopes at ranks seven and eight
are twelve and twenty. At $N=W+2=8$, the binary comparison (4)
is an equality: $(f_8-1)12-(f_7-1)20=20\cdot12-12\cdot20=0$.
The actual-reference all-shift C assertion is still valid; the proof
does not falsely promote this intermediate comparison to strictness.

\paragraph{Scope.}
The theorem proves unrestricted word/capacity/count/rank/shift C
in the ENTIRE homogeneous range, and uses written polynomial/Binet
identification with explicitly scoped Lean kernels and exact QA.
Earlier ranks $N\le W$, general actual C, general A2 and the
original unamplified one-descent unimodality conjecture remain OPEN.
No original least-counterexample numerical bound is upgraded.
